in such a way that the first factor permutes @@M0@@ and the second @@M1@@; @@M2@@ is embedded in the ring of invariants by sending @@M3@@ to the @@M4@@-th symmetric function of the @@M5@@; @@M6@@ sends the @@M7@@-th symmetric function of @@M8@@ (resp. @@M9@@) to the @@M10@@-th coefficient of @@M11@@ (resp. @@M12@@). Then:
@@M13@@
Since @@M14@@ is a homomorphism of groups, we are reduced to the case where @@M15@@, and in this case the additivity of @@M16@@ follows from the definition of the @@M17@@.

\sgaStableTarget{sga6.e05.para.1-7}1.7. We can generalize the preceding considerations to a product of functors @@M18@@; the functor @@M19@@ is representable by the ring
@@M20@@, @@M21@@, and @@M22@@, where @@M23@@ has weight @@M24@@. We denote by @@M25@@ the series @@M26@@; the element @@M27@@ of @@M28@@ has the universal property analogous to that of @@M29@@ (cf. 1.2).

We can introduce the morphisms @@M30@@, @@M31@@ as before, by repeating the construction for each family @@M32@@ separately. We obtain a result analogous to Theorem 1.6 by seeking the homomorphisms of functors from @@M33@@ to @@M34@@ that are @@M35@@-multilinear and compatible with the filtrations in each of the @@M36@@ arguments separately: they are determined by their value on the element
